\section{Limits}
\label{sec:limits}

\paragraph{\rGoPer{} bytes is a floor and will be quoted as a forecast.}
This is the most likely misuse of the paper. The measured goroutine holds a
place and nothing else. A working agent holds a conversation, a tool schema, and
whatever it has retrieved --- kilobytes to megabytes --- and that term dominates
\rGoPer{}\,bytes by three to four orders of magnitude. The correct reading is
narrow and still useful: \emph{the platform's own per-agent overhead is
negligible}, so a fleet's memory is the agents' data and not the runtime's
bookkeeping. Any capacity claim that multiplies \rGoPer{}\,bytes by a population
is wrong, including if we make it.

\paragraph{One host, one architecture.}
Every number is from one Apple M-series laptop. The goroutine and SQLite results
should port to any 64-bit platform with the same Go and SQLite versions --- they
measure allocator and file-format behaviour, not silicon --- but the container
figures will not: they were taken through a Linux VM on macOS, which adds a
virtualization layer that a native Linux host does not have. \rOci{}\,ms is
therefore likely \emph{pessimistic} for our own cold-start row. We have not
measured it natively and do not claim the better number.

\paragraph{History is the term that grows.}
A dormant agent stays \rFleetPer{}\,bytes forever. Its transcript does not. At a
retained 2\,KiB per call, a fleet serving a billion calls a month accumulates
roughly 22\,TB over twelve months \modeled{} --- which eventually exceeds
compute, and is the one line in the model with no ceiling. A platform built on
this decomposition must have a retention policy before it needs one; the cheap
dormant agent is what makes the expensive transcript easy to overlook.

\paragraph{Resume assumes the state is reachable.}
\rFleetResume{}\,ms is an in-process read from a local file. An agent whose state
lives on another node pays a network round trip instead, and the design that
makes this rare --- pinning an agent's tenant to a writer, as the companion
paper describes~\cite{unifiedtenant} --- has its own failure modes on failover.
The measurement is of the storage layer, not of a distributed system.

\paragraph{The in-process runtimes are not a boundary against a hostile tenant.}
Stated in \S\ref{sec:code} and repeated because it is the one place where taking
the cheap rung would be a security error rather than an economy: goja and
gpython run in the host address space. Where the tenant is adversarial, the third
rung is the only correct one, and its price is the price.

\paragraph{What we did not measure.}
Scheduler behaviour under contention at a million runnable (rather than parked)
goroutines; garbage-collection pause time at that heap; SQLite behaviour under
concurrent writers across many tenants; and the end-to-end latency of a real
agent turn, which is dominated by the model call and would tell you nothing about
the primitives this paper is about.
